\Needspace{10\baselineskip}
\subsection{\rev{Research directions}}
\label{sec:research_directions}
\begin{revision}
Future studies should establish which video measures can be used reliably in clinics, rehabilitation and home monitoring. The required evidence depends on the outcome and the decision it is intended to support.

\begin{itemize}
    \item \textbf{Clinically representative datasets.}
    Recruit participants across hip, knee and ankle conditions, severity levels, ages and body types. Record symptoms, walking speed, treatment history and compensation alongside synchronised reference measurements \cite{kaufman2001gait,weidow2006hip,ro2019effects,bacon2022relation}. Hold out participants and, where possible, clinical sites when assessing generalisation.

    \item \textbf{Biomechanical interpretability.}
    Connect a model's output to observable quantities such as joint-angle trajectories, asymmetry, range of motion or stance and swing timing \cite{ornetti2010gait,andriacchi2006role,astephen2008biomechanical,favre2016gait}. Specify the motion plane and gait phase and report how changes in walking speed or compensation affect the interpretation. Validate loading estimates against an appropriate biomechanical reference.

    \item \textbf{Multimodal assessment.}
    Pair video with inertial signals, muscle activity or clinical measures when they can add information about the outcome \cite{juri2020characterization,zhang2025myoelectric,zhang2024fall}. Compare the individual inputs and their combination on the same participants and task. That comparison can show whether any gain justifies the additional equipment and recording effort.

    \item \textbf{Robustness in everyday settings.}
    Record the same walks across relevant viewpoints, clothing, illumination and occlusion conditions \cite{ong2017efficacy,viswakumar2022omgait,van2023accuracy,wade2023examination}. Report error by joint and gait phase as well as the proportion of recordings that require correction or repetition. These measures are especially important when a method is intended for unsupervised home use.

    \item \textbf{Longitudinal and outcome-oriented validation.}
    Use repeat recordings to quantify measurement error and test whether video can detect functional changes during rehabilitation or disease progression \cite{gaudreault2011effects,bkaczkowicz2018gait,stenum2024clinical}. Relate observed changes to an appropriate clinical measure and evaluate whether they inform follow-up decisions.

    \item \textbf{Lightweight models for routine use.}
    Test portable systems on the devices and in the settings for which they are intended \cite{li2020human,uccheddu2022rgb,yan2025efficient}. Alongside runtime and acquisition effort, report agreement for the target joint variables. Efficient landmark estimation is most useful when the resulting gait measures retain their accuracy.

    \item \textbf{Clinical integration.}
    Work with clinicians and patients to decide which measurements fit a musculoskeletal assessment and how uncertainty should be reported. Pain, prior treatment and comorbidities can affect the interpretation of the same movement pattern \cite{hertling2006management,lu2026musculoskeletal,pandey2022movement}. Workflow evaluations should examine whether the output supports a decision and how much effort it takes to obtain.
\end{itemize}

These directions connect technical development to the practical question of when a gait measure is reliable enough to guide assessment. They also provide a basis for comparing systems tested in different settings.
\end{revision}
